\begin{exercice}\label{exoInf5} (\minsyndical)

On s'int\'eresse \`a un \'echantillon $(x_1,\ldots,x_n)$ dont la loi m\`ere suit une loi normale d'esp\'erance $\mu$ (\`a d\'eterminer) et de variance $2$.

\medskip

On veut d\'ecider entre les deux hypoth\`eses suivantes :

\begin{itemize}

\item $H_0$ : $X$ suit une loi normale d'esp\'erance $\mu_1 = 1$.

%vs.

\item $H_1$ : $X$ suit une loi normale d'esp\'erance $\mu_2 = 2$.

\end{itemize}

\medskip

On choisit comme statistique de test la moyenne de l'\'echantillon : $\bar{X} = \frac1n \sum_{i=1}^n X_i$.

\begin{enumerate}

\item On dispose d'un \'echantillon de taille $n=10$. On choisit comme intervalle d'acceptation $A = ] -\infty ; 1 + \delta ]$. Trouver $\delta$ pour que le risque $\alpha$ soit de $0.01$ (on souhaite que le risque soit tr\`es faible).

\item Calculer alors le risque de deuxi\`eme esp\`ece $\beta$ et la puissance du test $1-\beta$. 

\item M\^emes questions que 1. et 2. si on change le risque $\alpha = 0.05$.

\item M\^emes questions que 1. et 2. si on prend un \'echantillon plus grand : $n=500$ (et en gardant un risque $\alpha$ de 0.01).

\end{enumerate}

%\corrref{TD1_1}

\end{exercice}
